\documentclass[conference]{IEEEtran}
\IEEEoverridecommandlockouts

\usepackage{amsmath}
\usepackage{booktabs}
\usepackage{cite}
\usepackage{graphicx}
\usepackage{url}

\begin{document}

\title{MeshEcho: Confidence-Aware Routing for Airtime-Constrained LoRa Mesh IoT}

\author{
\IEEEauthorblockN{Anonymous Authors}
\IEEEauthorblockA{Submitted to IEEE ICC 2027, Internet of Things and Sensor Networks}
}

\maketitle

\begin{abstract}
LoRa mesh IoT networks trade delivery reliability for scarce airtime. Flooding
adds redundancy but creates collisions, whereas cached source routes are
efficient but can select a weak multi-hop path. We present MeshEcho, a
confidence-aware route-selection and bounded-recovery prototype, and evaluate
it with a reproducible packet-level simulator. The ICC release (version
2.1.15) uses a shared topology, traffic trace, independent reception stream,
and matched 2-s route-discovery window. In a calibrated 50-node, 8.25-km
multi-hop matrix over 20 seeds, MeshEcho obtains 0.720 ACK-confirmed PDR with
28.1 s aggregate airtime, compared with 0.537 and 38.5 s for the matched
source-route baseline. However, the Smart-CALM controller reaches 0.563 PDR,
and its no-confidence control reaches 0.382; the paired improvement is
separable from the controller's recovery behavior. A 3-seed topology gate
confirms non-saturated direct links and at least two graph hops. A legacy
contention case and an 18-km discovery stress test expose the limits of the
claim. The evidence supports a bounded reliability--airtime tradeoff, not
topology-independent dominance.
\end{abstract}

\begin{IEEEkeywords}
LoRa mesh, Internet of Things, sensor networks, adaptive routing, airtime
efficiency, packet delivery ratio
\end{IEEEkeywords}

\section{Introduction}
LoRa is attractive for low-power IoT because it combines long link budget with
small, inexpensive radios \cite{augustin2016study}. A single-hop star topology
is insufficient for obstructed or infrastructure-poor deployments, motivating
multi-hop and mesh extensions \cite{wong2024multihop,cotrim2020lorawanmesh,
centelles2021beyond}. Multi-hop forwarding is nevertheless expensive: a low
data rate makes every relay occupy the channel for a long time, while a flood
multiplies collision opportunities. A cached source route reduces redundancy,
but a route with one weak hop can fail even when another route exists.

MeshEcho explores the middle of this design space. It records lightweight
link-quality evidence during discovery, ranks candidate paths by confidence,
and invokes bounded recovery forwarding when discovery or a cached path is
uncertain. The implementation is intentionally modest: it is a simulation
prototype, not a replacement for a production mesh stack. The central
question is whether confidence can improve the reliability--airtime operating
point without being mistaken for a universal routing advantage.

This paper makes four contributions:
\begin{itemize}
\item a reproducible packet-level LoRa mesh harness with time-on-air,
propagation, shadowing, half-duplex collision/capture, route discovery, and
energy accounting;
\item a pre-specified, non-saturated multi-hop ICC matrix with per-seed quality
gates and matched discovery timing;
\item paired ablations that separate confidence ranking from timeout and
route-miss recovery; and
\item an explicit evidence boundary covering fixed-pair contention, route
discovery stress, and simulation-only limitations.
\end{itemize}

\section{Protocol and Model}
\subsection{Network model}
We simulate $N$ nodes uniformly placed in a square and sharing one half-duplex
LoRa channel. Unless stated otherwise, packets have 32-byte payloads, 17 dBm
transmit power, 125-kHz bandwidth, coding rate 4/5, capture threshold 6 dB,
and spreading factor SF7 for the calibrated matrix. Time-on-air follows the
standard LoRa symbol calculation \cite{semtech2013lora}; airtime is the sum of
all transmitted packet durations. Log-distance path loss with log-normal
shadowing produces received SNR. Reception probability is a logistic function
of SNR margin, and overlapping packets are accepted only when the desired
signal passes the capture rule. The simulator also records collision failures,
route-cache events, discovery success, and energy from configured TX/RX
currents.

\subsection{Routing policies}
Managed flooding (Meshtastic-like) rebroadcasts with duplicate suppression.
The source-route baseline (MeshCore-like) floods a request, caches a returned
path for a finite TTL, and forwards later data along that path. MeshCore-like
uses an immediate reply in the historical mode; ICC experiments use an
explicit 2-s candidate-collection window so reply timing is matched to CALM.
\texttt{calm-mesh} is the static MeshEcho policy: it ranks discovered routes
using path confidence and uses bounded route-miss fallback. \texttt{smart-calm}
adds an online profile controller; \texttt{smart-calm-no-fallback} disables
only route-miss fallback, and \texttt{smart-calm-no-confidence} disables
confidence-dependent admission and rewards. These controls are necessary
because a full-versus-baseline difference otherwise conflates mechanisms.

\section{Experimental Design}
\subsection{Calibrated ICC matrix}
The primary matrix has 50 nodes in an $8250\,\mathrm{m}$ square, 600 s per
run, mixed traffic at one application flow per minute, SF7, path-loss exponent
2.75, 4-dB shadowing, and 20 seeds. For each seed, 24 source--destination
pairs are selected from a shared analytical link-budget graph: retained edges
have PRR at least 0.90 and selected graph distance is 2--4 hops. The quality
gate rejects a seed unless all 24 pairs are present, mean selected distance is
at least 2.0 hops, and at least 10\% of direct links have PRR below 0.99.
Reception draws use an independent random stream from protocol jitter and
exploration, preventing later protocols from inheriting a changed channel
trace. Smart-CALM timeout retries are zero. Values below are means over 20
seeds with two-sided 95\% confidence half-widths.

\begin{table*}[t]
\centering
\caption{Primary calibrated multi-hop matrix (20 seeds). ACK PDR is
ACK-confirmed unicast delivery; route fraction is among completed cached
routes.}
\label{tab:main}
\scriptsize
\setlength{\tabcolsep}{3.5pt}
\begin{tabular}{lrrrrrr}
\toprule
Protocol & ACK PDR & Dest. PDR & Airtime (s) & Discovery & Multi-hop route & Mean cached hops\\
\midrule
MeshCore-like & .537$\pm$.081 & .679$\pm$.093 & 38.5$\pm$5.3 & .741 & .775 & 1.80\\
MeshEcho (CALM) & \textbf{.720$\pm$.115} & \textbf{.739$\pm$.115} & \textbf{28.1$\pm$4.7} & .891 & .989 & 2.03\\
Smart-CALM & .563$\pm$.088 & .630$\pm$.100 & 28.6$\pm$4.7 & .801 & 1.000 & 2.13\\
Smart-CALM, no fallback & .561$\pm$.081 & .607$\pm$.096 & 28.0$\pm$4.5 & .811 & 1.000 & 2.11\\
Smart-CALM, no confidence & .382$\pm$.122 & .470$\pm$.133 & 27.8$\pm$4.6 & .589 & .730 & 1.61\\
\bottomrule
\end{tabular}
\end{table*}

\subsection{Controls and robustness cases}
The same release reruns a legacy 3000-m repeated-pair matrix (50 nodes,
1200 s, six flows/min, SF9) for traceability, plus a 6-dB shadowing case and
a ten-flows/min load case. These are contention-oriented sensitivity cases,
not substitutes for the calibrated matrix. A separate 18000-m, 3-seed probe
uses the same connected-pair selector with a PRR threshold of 0.70. It is a
topology and route-discovery quality gate, not a headline PDR comparison.

\section{Results}
\subsection{Calibrated multi-hop behavior}
Table~\ref{tab:main} establishes that the primary comparison is genuinely
multi-hop rather than a direct-link replay. Every seed passed the gate; the
mean selected graph distance is exactly 2.0 hops, while 16.0\% of direct links
are below 0.99 PRR and 4.6\% are below 0.50. MeshEcho improves over the
matched source-route baseline by 0.183 ACK PDR and reduces airtime by 26.8\%.
The result is a protocol-bundle comparison: CALM's bounded recovery and
confidence ranking are both enabled.

The ablations make this boundary measurable. Smart-CALM minus
no-confidence is +0.181 ACK PDR with a paired 95\% interval
$[+0.094,+0.268]$, showing a confidence-sensitive difference in this
calibrated regime. Smart-CALM minus no-fallback is only +0.002 with interval
$[-0.020,+0.024]$; its mean PDR gain is therefore not isolated from fallback
in this replay. The no-confidence line also has lower discovery success and
fewer multi-hop cached routes, so it is evidence about the complete disabled
policy, not a controlled oracle with identical route candidates.

\begin{table}[t]
\centering
\caption{2.1.15 legacy mixed-traffic sensitivity case (20 seeds). This table
is retained for traceability; direct links are mostly easy and paths are not a
neutral deployment benchmark.}
\label{tab:legacy}
\scriptsize
\setlength{\tabcolsep}{2.5pt}
\begin{tabular}{lrrr}
\toprule
Protocol & ACK PDR & Broadcast & Airtime (s)\\
\midrule
Managed flooding & .853 & .942 & 1177\\
MeshCore-like & .650 & .958 & 898\\
MeshEcho (CALM) & .728 & .970 & 847\\
Smart-CALM & .936 & .962 & 827\\
\bottomrule
\end{tabular}
\end{table}

\subsection{Legacy robustness and discovery stress}
In the repeated-pair mixed case (Table~\ref{tab:legacy}), MeshEcho saves 5.7\%
airtime versus MeshCore-like while raising ACK PDR by 0.078. The high-load
case reduces all delivery ratios, while the 6-dB shadowing case preserves the
same ordering within the uncertainty of the frozen workload. These cases are
useful for regression and contention behavior, but repeated pairs favor route
reuse and the direct-link PRR distribution is saturated.

The 18-km quality probe passed all three seeds with mean graph distance 2.014,
direct-link fractions below 0.99 and 0.50 of 0.641 and 0.445, and a complete
24-pair pool. MeshCore-like discovery succeeded on only 0.178 of attempts in
this deliberately harsh setting. We therefore report it as a route-discovery
stress diagnostic; treating its low PDR as a general source-route result would
confound discovery timing, connectivity, and data forwarding.

\section{Fairness, Limitations, and Reproducibility}
All protocols in a run share node placement, static shadowing, traffic pairs,
PHY parameters, seed, and reception stream. The primary matrix also matches
the MeshCore-like discovery window and route-cache lifetime. Nevertheless,
the controls do not make every native budget identical: bounded fallback,
profile switching, and route-repair semantics differ by design. The calibrated
matrix is therefore a fair shared-harness protocol-bundle comparison, while
the paired ablations provide narrower mechanism evidence.

This is simulation-only evidence. The model omits waveform-level interference,
duty-cycle enforcement, queue implementation details, time-varying fading,
variable route-header size, and hardware sensitivity variation. The confidence
estimate uses the same static link realization as forwarding, so stale-quality
behavior is not tested. The 3-seed stress gate is a quality check, not a power
analysis. Future work should preregister a larger pair-churn matrix, equalize
retry and fallback budgets across baselines, model variable headers, and
validate the policy on hardware or hardware-in-the-loop.

All source, tests, experiment scripts, raw CSVs, reports, and this paper are
released in version 2.1.15 at
\url{https://github.com/BH4ME/lora-mesh-routing-sim}. The main workflow is
\texttt{OUT\_PREFIX=meshecho\_v2\_1\_15\_icc2027 bash tools/run\_icc\_experiments.sh};
the script runs the four legacy cases, the topology gate, and the calibrated
matrix. No human participants or personal data were involved; there are no
known conflicts of interest or external funding.

\section{Conclusion}
MeshEcho reaches a useful reliability--airtime operating point in a calibrated
multi-hop LoRa simulation: 0.720 ACK PDR and 28.1 s airtime versus 0.537 and
38.5 s for a matched source-route baseline. The result is reproducible and
passes explicit topology/link-quality gates. However, paired controls show
that Smart-CALM's confidence-sensitive difference and its recovery behavior
must be reported separately, while the harsh discovery probe warns against
using route-discovery failure as a universal protocol ranking. Version 2.1.15
therefore supports a bounded mechanism and tradeoff claim, not a claim of
general dominance or physical validation.

\bibliographystyle{IEEEtran}
\bibliography{references}

\end{document}
